\section{Introduction}

Personalized text-to-image generation has progressed rapidly in recent years, especially in the single-subject regime where a model must preserve one reference identity while following a textual prompt. However, practical deployment increasingly requires multiple distinct reference identities to co-exist in one image, often with explicit interactions, physical contact, asymmetric role assignments, or object bindings. In this regime, the problem is no longer just recognition or prompt following. The real bottleneck is \emph{binding}: the model must preserve who is who, which appearance belongs to whom, and which interaction is assigned to which subject.

The evaluation layer has not kept pace with this difficulty. In many current pipelines, quality is still measured with global similarity metrics (e.g., CLIP \cite{radford2021learning}) or general preference models (e.g., ImageReward \cite{xu2023imagereward}). These tools are useful for coarse semantic plausibility, but they are poorly matched to the real errors that dominate multi-subject personalization. A generated image can mention the right scene, contain approximately the right colors, and still be unusable because one subject is missing, two identities collapse into generic clones, clothing attributes bleed across people, or a requested interaction is assigned to the wrong subject. In such cases, a useful evaluator must preserve ranking separability even when the number of requested subjects becomes large.

This paper is motivated by a simple empirical observation: once the subject count increases, existing metrics often lose the ability to separate obviously better generations from obviously worse ones. If a benchmark cannot distinguish ``missing one subject'' from ``missing most subjects,'' then it cannot support real model development or credible ablation. In other words, binding is not only the bottleneck for generation; it is also the bottleneck for evaluation. Constructing a high-quality, human-annotated dataset for multi-subject personalization at scale is prohibitively expensive, which has stalled the development of better evaluation tools.

To overcome this, we introduce \textbf{BindBench}, a large-scale, rigorously controlled benchmark and dataset for multi-subject generation. BindBench is designed around a decoupled paradigm that balances scale and rigor:
\begin{itemize}
    \item \textbf{BindBench-Silver (60K pairs):} A massive training ground generated by internal models (Nano/Mosaic) and labeled by dual teacher VLMs (Gemini 2.5 Flash and Gemini 3.1 Flash Lite). The prompts are generated via a strict hierarchical construction, spanning $N \in \{2,4,6,8\}$ entities across 36 controlled buckets of interaction types and human/object ratios.
    \item \textbf{BindBench-Gold (4K pairs):} A carefully curated, double-blind human-annotated evaluation set covering six diverse state-of-the-art generators, including unseen models like Flux, Seedream 4.5, and GLM. This serves as the ultimate cross-generator meta-evaluation ground.
\end{itemize}

By decoupling the data, BindBench provides both an abundant, low-cost training bed for automated metrics and an unbiased testbed for evaluating generation quality. 

To demonstrate the utility of this benchmark, we propose \textbf{AnchorScore}, a learned evaluator trained exclusively on the 60K silver set. AnchorScore takes the references, prompt, and generated image pair as input, then predicts both pairwise preference and three interpretable image-level diagnostic dimensions: \emph{Existence}, \emph{Appearance}, and \emph{Interaction}. Our experiments reveal that AnchorScore exhibits strong zero-shot generalization on the 4K gold set, outperforming standard metrics on unseen generators. This validates our core thesis: the highly structured 60K silver set is a powerful asset not just for training robust evaluators, but potentially for aligning generative models (e.g., via DPO/RLHF) to solve multi-subject entanglement.

In summary, this paper makes three contributions:
\begin{itemize}
    \item We introduce \textbf{BindBench}, a controlled benchmark with hierarchical prompt construction, factorized scene buckets, a 60K silver training resource, and a 4K cross-generator human evaluation set.
    \item We provide \textbf{AnchorScore}, a strong reference-conditioned evaluator baseline that couples pairwise ranking with interpretable diagnostic prediction.
    \item We establish a benchmark protocol that simultaneously measures generator quality, metric failure, and cross-generator generalization in high-subject-count personalized generation.
\end{itemize}

The expected outcome of this work is a comprehensive benchmark-and-metric stack. Section~2 reviews existing evaluation paradigms. Section~3 defines the problem setup. Section~4 introduces the hierarchical construction of BindBench. Section~5 presents the AnchorScore baseline evaluator. Section~6 describes our compute-controlled fair ablation protocol. Section~7 outlines the experiments benchmarking both the generators and the metrics, and Section~8 concludes with future work.
